% Standalone research entry 217; batch 208--217.
% Original section material preserved from Pasted text(3).txt.
% Research additions: 27 September 2026.
% Compile twice with: lualatex 217.tex
% !TeX program = lualatex
% Compile twice: lualatex open_problems_500.tex
% All references and prose are embedded; no BibTeX database or external inputs.
\documentclass[11pt,a4paper]{article}
\usepackage[margin=24mm,headheight=14pt]{geometry}
\usepackage{fontspec}
\setmainfont{Latin Modern Roman}
\setsansfont{Latin Modern Sans}
\setmonofont{Latin Modern Mono}
\usepackage{amsmath,amssymb,mathtools}
\usepackage{unicode-math}
\setmathfont{Latin Modern Math}
\usepackage{microtype}
\usepackage{enumitem,xurl,needspace}
\usepackage[dvipsnames]{xcolor}
\usepackage{hyperref}
\hypersetup{unicode=true,colorlinks=true,linkcolor=MidnightBlue,urlcolor=MidnightBlue,
 pdftitle={500 Mathematical Problem Statements},pdfauthor={Source-annotated compilation},bookmarksnumbered=true}
\usepackage{bookmark}
\usepackage{tocloft}
\setlength{\cftsecnumwidth}{3em}
\cftsetpnumwidth{2.5em}
\cftsetrmarg{4em}
\usepackage{titlesec}
\titleformat{\section}{\Large\bfseries\raggedright}{\thesection}{0.7em}{}
\usepackage{fancyhdr}
\pagestyle{fancy}\fancyhf{}
\fancyhead[L]{\small\sffamily Mathematical problem statements}
\fancyhead[R]{\small Source edition 22}
\fancyfoot[C]{\thepage}
\setlength{\parindent}{0pt}
\setlength{\parskip}{5pt plus 1pt}
\setlength{\emergencystretch}{3em}
\setcounter{tocdepth}{1}
\setcounter{secnumdepth}{1}
\setlist[itemize]{leftmargin=3em,itemsep=5pt,topsep=4pt,parsep=0pt}
\allowdisplaybreaks
\newcommand{\N}{\mathbb{N}}
\newcommand{\Z}{\mathbb{Z}}
\newcommand{\Q}{\mathbb{Q}}
\newcommand{\R}{\mathbb{R}}
\newcommand{\C}{\mathbb{C}}
\newcommand{\F}{\mathbb{F}}
\newcommand{\A}{\mathbb{A}}
\newcommand{\PP}{\mathbb P}
\newcommand{\E}{\mathbb{E}}
\newcommand{\eps}{\varepsilon}
\newcommand{\dd}{\,\mathrm d}
\newcommand{\sref}[2]{\hyperlink{source:#1:#2}{[S#2]}}
\newcommand{\eref}[2]{\hyperlink{extra:#1:#2}{[E#2]}}
% Turn the next switch off for a shorter reading copy. The quotations preserve
% every exactTarget field, including qualifications omitted from a short summary.
\newif\ifshowcatalogue
\showcataloguetrue
\newcommand{\cataloguescope}[1]{\ifshowcatalogue
 \paragraph{Exact scope in the supplied catalogue.}
 \begin{quote}\small #1\end{quote}\fi}

% Standalone wrapper; no external inputs or bibliography files are required.
\fancyhead[L]{\small\sffamily Research notebook: section 217}
\fancyhead[R]{\small 27 September 2026}
\hypersetup{pdftitle={Research attempt 217},pdfauthor={Research notebook}}
\setlength{\parskip}{4pt plus 1pt}
\setlist[itemize]{before=\small\interlinepenalty10000,itemsep=3pt}
\begin{document}
\setcounter{section}{216}
% ================================================================
% problemId: problem.skolem-problem-for-linear-recurrence-sequences
% source releaseRank: 217; source releaseStatus: open_with_solved_subcases
% exactTarget SHA-256: 5f6760ab4fb6873bb3392f1a15a17ccffc8b54728753f6b086de62054a1b0e00
\clearpage
\section{\texorpdfstring{Skolem Problem for Linear Recurrence Sequences}{Skolem Problem for Linear Recurrence Sequences}}\label{217}
\subsection{Definitions and mathematical statement}
A linear recurrence sequence of order $k$ is determined by coefficients $c_0,\ldots,c_{k-1}$ and initial terms $u_0,\ldots,u_{k-1}$ through
\[
 u_{n+k}=c_{k-1}u_{n+k-1}+\cdots+c_0u_n\quad(n\ge0).
\]
The data are integers or algebraic numbers given by effective exact encodings. The Skolem problem asks for an algorithm that always halts and decides whether $\exists n\ge0$ with $u_n=0$. The order is part of the input and is unbounded. An algorithm for a fixed low order, or a description of the zero set without effective bounds on the exceptional indices, does not settle the general decision problem.
\par\smallskip\textit{Formulation and concept references:} \sref{217}{1}.
\cataloguescope{Decide algorithmically, given an integer (or algebraic) linear recurrence sequence, whether it has a zero term.}
\subsection{Short English statement}
Can one always decide from a linear recurrence and its initial values whether any term of the resulting sequence is exactly zero?
% BEGIN RESEARCH ADDITION 217
\subsection{Research attempt}

\paragraph{Current frontier.}
Bacik's order-four theorem includes algebraic coefficients and initial
values \sref{217}{1}; the order is not bounded by four in this entry.
The Skolem--Mahler--Lech description gives a finite union of arithmetic
progressions together with a finite exceptional zero set, but finiteness
alone does not bound the last exceptional index. Luca--Ouaknine--Worrell's
July 2026 manuscript \eref{217}{1}, Conjecture 4.2 and Theorem 4.3, obtains
a conditional large-zero bound from a spacing conjecture for a specified
set of ``good primes.'' That conjecture is stronger than invoking the
usual Cram\'er--Granville conjecture for all primes. The manuscript's
unconditional sparsity statements do not decide whether an individual
exceptional index exists. These distinctions leave room for both
archimedean estimates and local certificates, examined below without
assuming the prime-gap conjecture.

\paragraph{Attack route.}
Enumerating terms semidecides the existence of a zero. To make a total
decision procedure, seek either a finite modular certificate of zero-freeness,
a $p$-adic zero isolation argument that rules out integer indices, or an
effective dominating-term estimate. The modular strategy is attractive
because every finite residue orbit is computable. I actively test its
completeness and find an exact order-seven counterexample to that
\emph{intermediate strategy}, not to decidability. The dominating-term
route then supplies a genuine terminating subcase and isolates what
fails when several characteristic roots have the same modulus.

\paragraph{Attempt.}
\textit{1. Finite residue tests are sound but not complete.}
For an integer recurrence of order $k$, the state vector of $k$ successive
terms evolves deterministically on $(\Z/m\Z)^k$. By recording states
until repetition, one can decide whether any term is zero modulo $m$.
If none is, then the integer sequence has no zero. It is tempting to run
this search over $m=2,3,\ldots$ in parallel with exact zero enumeration,
expecting one side to halt. The following sequence disproves that
expectation:
\[
 u_n=(n^2-13)(n^2-17)(n^2-221),\qquad n\ge0.                \tag{217.1}
\]
It never vanishes at an integer, since $13,17,221$ are not integer
squares. As a polynomial in $n$ of degree six, it satisfies
$(T-1)^7u=0$, where $Tu_n=u_{n+1}$. Explicitly,
\[
 \begin{split}
 u_{n+7}={}&7u_{n+6}-21u_{n+5}+35u_{n+4}-35u_{n+3}\\
          &+21u_{n+2}-7u_{n+1}+u_n.                         \end{split}\tag{217.2}
\]
Thus it is a valid input of order seven to the original problem.

Nevertheless, $u_n\equiv0\pmod m$ is solvable for every $m\ge1$.
For an odd prime $p\notin\{13,17\}$, if neither $13$ nor $17$ is a
quadratic residue, their product $221$ is. Choose a square root modulo
$p$ of one of these nonzero residues; its derivative $2n$ is nonzero,
so the root lifts to every $p$-power. At $p=13$, $17\equiv2^2$;
at $p=17$, $13\equiv8^2$. These roots also lift. At $p=2$,
$17\equiv1\pmod8$ has a square root in $\Z_2$, hence modulo every
power of two. The last elementary fact follows by lifting an odd square
root one binary digit at a time: replacing an odd root modulo $2^j$
by that root plus $2^{j-1}$ toggles its square modulo $2^{j+1}$.
The Chinese remainder theorem combines the chosen prime-power roots,
possibly from different factors, into one nonnegative index modulo $m$.
This proves the asserted failure for \emph{all} moduli, not just those
searched by a program.

The reproduction script verifies (217.2) symbolically and finds a root
of (217.1) modulo every integer $2\le m\le256$. The finite check is
only a test of the implementation; the prime-power argument proves the
unbounded statement. The polynomial nature of (217.1) makes its actual
zero-freeness easy to decide. The example therefore does not imply that
Skolem is undecidable; it proves incompleteness of this particular modular
certificate scheme. Its only characteristic root is $1$ with multiplicity
seven, so the obstruction also highlights the distinction between
``simple'' characteristic roots and the usual nondegeneracy condition
on quotients of \emph{distinct} roots.

\textit{2. What $p$-adic interpolation does and does not supply.}
For clarity, first consider an integer matrix representation
$u_n=\ell A^n v$ with $\det A\ne0$. Choose an odd prime $p$ not dividing
$\det A$. Some computable $h>0$ satisfies $A^h\equiv I\pmod p$,
so write $A^h=I+pB$ with $B$ integral. For each $0\le r<h$, the function
\[
 F_r(z)=\ell A^r\sum_{j=0}^{\infty}
                     \binom zj p^jB^jv,\qquad z\in\Z_p,    \tag{217.3}
\]
interpolates $u_{r+hn}$ at every nonnegative integer $n$.
The convergence is uniform because $\binom zj\in\Z_p$ for $z\in\Z_p$.
Moreover the coefficients of the polynomial $\binom zj$ have
$p$-adic sizes at most $p^{v_p(j!)}$, and
$j-v_p(j!)\to\infty$ for odd $p$. Thus the series converges in the
algebra of restricted power series on $\Z_p$: (217.3) is analytic,
not just continuous. A nonzero such series has finitely many zeros by
Strassmann's theorem. The standard interpolation perspective underlies
the structural zero-set theorem discussed in \eref{217}{1}, Section 1.

The missing step is arithmetic, not the uniform convergence just proved.
A zero of $F_r$ may be a $p$-adic integer without being an ordinary
nonnegative integer. Finitely many isolated $p$-adic roots do not by
themselves give a terminating procedure for recognizing that none is an
ordinary integer. Increasing $p$-adic precision only narrows congruence
classes, each of which contains arbitrarily large nonnegative integers.
In (217.1), the interpolating polynomial has a local root for every prime
but no integer root. Algebraicity of its roots resolves that particular
example; it is not established for roots of every analytic interpolant
(217.3). Thus an integer-recognition theorem would be a substantive extra
input, not an automatic application of zero isolation.

\textit{3. A completely effective dominant-root subcase.}
There is a useful escape from this obstruction when one archimedean term
dominates. Effective algebraic linear algebra gives, after any finite
initial transient coming from the root zero, an expression
\[
 u_n=\sum_{i=0}^s P_i(n)\lambda_i^n                       \tag{217.4}
\]
with distinct nonzero algebraic $\lambda_i$ and algebraic-coefficient
polynomials $P_i$. Terms with zero polynomials are deleted. Fix one complex
embedding. Suppose
\[
 |\lambda_0|>\max_{i\ge1}|\lambda_i|,
 \qquad P_0\ne0.                                         \tag{217.5}
\]
All comparisons needed here can be decided for algebraic real numbers
$|\lambda_i|^2$. Let $r=\deg P_0$ and $a\ne0$ be its leading coefficient.
Effective rational bounds on its coefficients give a computable $N_0$
and a positive rational $c<|a|/2$ such that
\[
 |P_0(n)|\ge c n^r\quad(n\ge N_0).
\]
For example, for $r\ge1$, bound the lower terms by
$n^{r-1}\sum_{j<r}|a_j|$ and take
$n\ge 2\sum_{j<r}|a_j|/|a|$, with rational upper and lower bounds used
in the computation. For $r=0$ the estimate holds immediately.
Choose a computable rational $\rho$ with
$\max_{i\ge1}|\lambda_i/\lambda_0|<\rho<1$, and rational $C>0$ and
integer $s_0\ge0$ such that
\[
 \sum_{i\ge1}|P_i(n)|\le C n^{s_0}\quad(n\ge1).
\]
The case with no other terms is simpler and is handled by the leading
polynomial bound alone. Otherwise put $d_0=\max(s_0-r,0)$ and search for
an integer $N\ge\max(N_0,1)$ satisfying
\[
 C N^{d_0}\rho^N<c,
 \qquad \rho(1+1/N)^{d_0}<1.                              \tag{217.6}
\]
This is an exact rational search. It terminates because exponential decay
beats every fixed polynomial. The second inequality makes
$n^{d_0}\rho^n$ decreasing for all integers $n\ge N$. Therefore
\[
 \left|\sum_{i\ge1}P_i(n)(\lambda_i/\lambda_0)^n\right|
   \le C n^{s_0}\rho^n<c n^r\le |P_0(n)|\quad(n\ge N).
\]
So $u_n\ne0$ throughout the tail. Checking the finitely many earlier
indices by exact algebraic arithmetic decides existence of a zero.
This rederives a standard dominance mechanism with an explicit stopping
condition; it is not a claim of a newly solved unrestricted class.
It also treats the one-characteristic-root polynomial example (217.1).

\textit{4. The remaining oscillatory regime.}
When several $|\lambda_i|$ share the maximum, normalization of (217.4)
leaves a sum of oscillating algebraic phases with polynomial amplitudes.
The factor $\rho<1$ used in (217.6) is no longer available for their
mutual cancellation. A successful replacement needs an effective
separation estimate, or a different finite certificate, valid for every
nonzero value in the tail. A bound on the number of zeros is not such an
estimate: one zero could occur after every cutoff reached so far. Nor does
a density-one set of decidable indices cover the exceptional complement.
The conditional good-prime argument in \eref{217}{1}, Sections 4--6,
addresses precisely a large-index issue, but its spacing hypothesis is
not established by the present reasoning.

\paragraph{Stress test.}
The modular example uses only integer data and hence genuinely lies inside
the original algebraic-data scope. Its zero-free proof covers every
nonnegative index, and its local-solvability proof covers every modulus.
It refutes a proposed completeness lemma rather than the existence of all
possible algorithms. The interpolation discussion is explicitly restricted
to invertible integral $A$; no claim that every input immediately has that
form is needed for the identified obstruction. In the dominance argument,
zero characteristic roots and zero coefficient polynomials are removed
with their finite transients accounted for; the remaining finite-prefix
checks are exact, not floating-point root tests.

There is a methodological relation to Section 208: an ineffective finite
exceptional set does not supply a stopping rule. No reduction from Effective
Roth to full Skolem, or conversely, has been proved here. In particular,
neither the known order-four theorem \sref{217}{1} nor the conditional 2026
prime-gap theorem may be applied with its order or conjectural hypothesis
deleted.

\paragraph{Outcome.}
\textbf{Outcome: Exploratory approach with unresolved gap.}

The attempt gives a proved obstruction to completeness of all finite
modular zero-freeness tests, a carefully justified $p$-adic interpolation
with its integer-recognition gap exposed, and an explicit terminating
algorithm when a unique characteristic root has strictly maximal modulus.
These results do not decide arbitrary-order recurrences in the remaining
oscillatory regime. No general zero cutoff, unconditional replacement for
the prime-gap hypothesis, or counterexample to algorithmic decidability
has been obtained.
% END RESEARCH ADDITION 217
\subsection{Sources}
\begin{itemize}\raggedright
\item[\textnormal{[S1]}]\hypertarget{source:217:1}{} Piotr Bacik, "Completing the picture for the Skolem Problem on order-4 linear recurrence sequences," TheoretiCS, Volume 4 (published 2 December 2025), DOI 10.46298/theoretics.25.28; arXiv:2409.01221.
\url{https://theoretics.episciences.org/17020}.
{\small\textit{Catalogue formulation source.}}
% BEGIN RESEARCH REFERENCES 217
\item[\textnormal{[E1]}]\hypertarget{extra:217:1}{} F. Luca, J. Ouaknine and J. Worrell, ``Conjectural Decidability of the Skolem Problem,'' arXiv:2607.15510v1 (16 July 2026), Theorem 1.1, Conjecture 4.2, Theorem 4.3, and Sections 5--6. The good-prime spacing assumption is conjectural, and sparsity alone is not full decidability.
\url{https://arxiv.org/html/2607.15510v1}.
% END RESEARCH REFERENCES 217
\end{itemize}

\end{document}
